\section{低秩 Taylor 矩、有理包围与两个试点证书}\label{sec:certificates}

本节证明两个具体单元的预算严格小于 $1$。
计算的数学对象是式 \eqref{dens:B} 给出的完整预算，
包括无限高度剩余权；数值证书不替代前两节的解析推导。
我们先把双重有限和分解为非负矩，再将每个实数因子向上包围为
分母 $2^{48}$ 的有理数，最后仅用整数与有理数运算验证目标界。
生成器可以提出这些有理数，但每一个包围关系另有 Lean 证明。

\subsection{把完整零点预算精确拆成主项与辅助项}

沿用实部网格 $\sigma_0,\ldots,\sigma_J$、高度网格
$u_0,\ldots,u_M=U$、左端点 $A$ 和参数 $q$，并假定 $M\geq1$。
这一条件保证完整高度余项可以放入最后一个高度单元。
对 $0\leq i<M$ 记
\begin{align}
 s_i&=u_{i+1},&
 v_i&=s_i+\log s_i,&
 \mu_i&=g_q(\e^{u_i})-g_q(\e^{s_i}),\label{cert:height-data}\\
 d_i&=2s_i^2\mu_i+
  \mathbf{1}_{\{i=M-1\}}\,32c_2q^2\e^{-2s_i}s_i^2. \label{cert:d}
\end{align}
其中 $s_i>0$，$d_i\geq0$。最后一行的附加项
恰来自式 \eqref{dens:B} 中 $16$ 倍的完整余项，
并已乘上正、负高度合并产生的因子 $2$。
遗漏此项会得到不完整的证书。
再定义
\begin{equation}\label{cert:cDE}
 c_j=B_{1,r(j)}W_j(A),\qquad
 D_A=\sum_jB_{2,r(j)}W_j(A),\qquad
 E_A=\sum_jB_{3,r(j)}W_j(A).
\end{equation}
这里 $D_A,E_A$ 只是密度辅助系数，与
第 \ref{sec:trapezoid} 节的差分算子和平凡零点差分不是同一对象。

因为 $s>0$ 时
\[
 \e^{ps}s^{p+2}=s^2\e^{p(s+\log s)},
\]
式 \eqref{dens:Phi} 可改写为
\[
 \frac{\Phi_A(s)}{s^2}
  =\sum_jc_j\e^{p_j(s+\log s)}
          +D_A+\frac{E_A}{s}+\frac{\e^{-A/2+s}}{s}.
\]
由 $d_i$ 的定义逐项展开，可得精确恒等式
\begin{equation}\label{cert:decomposition}
 2\mathcal B(A,q;\boldsymbol u)=\mathcal M+\mathcal A,
\end{equation}
其中
\begin{align}
 \mathcal M&=\sum_{j=0}^{J-1}\sum_{i=0}^{M-1}
                   c_jd_i\e^{p_jv_i},\label{cert:main}\\
 \mathcal A&=\sum_{i=0}^{M-1}d_i
       \left(D_A+\frac{E_A}{s_i}+\frac{\e^{-A/2+s_i}}{s_i}\right).
                                                     \label{cert:aux}
\end{align}
此处没有数值近似，也没有放弃任何一项。
对应
\leanref{M2PilotBudgetIdentity}{weightedPhi_eq_twice_heightBudget}、
\leanref{M2PilotBudgetIdentity}{weightedPhi_eq_main_add_aux}
以及远段的同名恒等式。

\subsection{有余项的 Taylor 不等式}

\begin{lemma}\label{cert:taylor}
对 $0\leq z\leq11$，
\begin{equation}\label{cert:taylor-eq}
 \e^z\leq\sum_{\ell=0}^{20}\frac{z^\ell}{\ell!}
                         +2\frac{z^{21}}{21!}.
\end{equation}
\end{lemma}
\begin{proof}
指数级数中从第 $21$ 项开始，相邻项之比为
$z/(\ell+1)\leq11/22=1/2$。
因此这部分尾和不超过首项 $z^{21}/21!$ 的两倍。
前 $21$ 项保留不变，即得到结论；$z=0$ 也包含在内。
\end{proof}
对应 \leanref{LowRankTaylor}{exp_le_taylor20}。
这不是无余项截断多项式；系数 $2$ 专门吸收全部 Taylor 尾和。

选取移位常数 $p_0,v_0$，满足 $p_j\geq p_0$、$v_i\geq v_0$。
定义四组非负数
\begin{equation}\label{cert:shifted-factors}
 f_j=c_j\e^{p_jv_0},\quad a_j=p_j-p_0,\qquad
 g_i=d_i\e^{p_0(v_i-v_0)},\quad b_i=v_i-v_0 .
\end{equation}
恒等式
\[
 p_jv_i=p_jv_0+p_0(v_i-v_0)+(p_j-p_0)(v_i-v_0)
\]
给出
\begin{equation}\label{cert:factorization}
 \mathcal M=\sum_{j,i}f_jg_i\e^{a_jb_i}.
\end{equation}
若还满足 $a_jb_i\leq11$，则可应用引理 \ref{cert:taylor}。
对非负数列 $(f,a)$ 定义
\[
 M_\ell(f,a)=\sum_j f_ja_j^\ell .
\]
把有限双重和中的多项式逐项展开并交换求和，得到
\begin{equation}\label{cert:low-rank}
 \mathcal M\leq
 \mathcal L_{20}(f,a;g,b):=
 \sum_{\ell=0}^{20}
       \frac{M_\ell(f,a)M_\ell(g,b)}{\ell!}
       +2\,\frac{M_{21}(f,a)M_{21}(g,b)}{21!}.
\end{equation}
原来 $J M$ 个指数项被化为两组各 $22$ 个矩。
这一降维仅使用分配律和有严格余项的 Taylor 不等式；
没有用低秩数值拟合替换原函数。
对应 \leanref{LowRankTaylor}{shifted_exp_factorization}
和 \leanref{LowRankTaylor}{shifted_exp_sum_le}。

\subsection{逐因子向外包围与精确整数矩}

因为 \eqref{cert:low-rank} 中的系数全部非负，
如果四组因子分别向上放大，则右边不会减小。
具体选择非负整数
$F_j,A_j,G_i,B_i$，并逐项证明
\begin{equation}\label{cert:dyadic-bounds}
 0\leq f_j\leq\frac{F_j}{2^{48}},\quad
 0\leq a_j\leq\frac{A_j}{2^{48}},\qquad
 0\leq g_i\leq\frac{G_i}{2^{48}},\quad
 0\leq b_i\leq\frac{B_i}{2^{48}}.
\end{equation}
于是
\[
 M_\ell(f,a)\leq
 \frac{I_\ell}{2^{48(\ell+1)}},
 \quad I_\ell=\sum_j F_jA_j^\ell,
 \qquad
 M_\ell(g,b)\leq
 \frac{J_\ell}{2^{48(\ell+1)}},
 \quad J_\ell=\sum_i G_iB_i^\ell .
\]
从而得到完全有理的上界
\begin{equation}\label{cert:rational}
 \mathcal R=
 \sum_{\ell=0}^{20}
  \frac{I_\ell J_\ell}{\ell!\,2^{96(\ell+1)}}
  +\frac{2I_{21}J_{21}}{21!\,2^{96\cdot22}},
 \qquad \mathcal M\leq\mathcal R.
\end{equation}
欲证明 $\mathcal R\leq r$，只须乘上一个公共正整数分母，
验证一个整数不等式。有限有理运算不会产生浮点舍入误差。

必须区分三个可分别审计的步骤。首先，实数表达式
$f_j,a_j,g_i,b_i$ 由 \eqref{cert:height-data}--\eqref{cert:shifted-factors}
精确定义。其次，\eqref{cert:dyadic-bounds} 是需要证明的不等式，
不能因外部程序输出了某个数就认定它成立。最后，
\eqref{cert:rational} 对已经固定的整数表做精确算术。
本实现的生成脚本负责提出表和包围关系；Lean 的区间引理证明第二步，
内核归约验证第三步。对应
\leanref{LowRankTaylor}{lowRankTaylor20_mono}、
\leanref{LowRankInteger}{lowRankTaylor20_eq_rational}。

具体的有向包围遵守如下规则。例如，若 $s>0$ 且
$\e^\ell\leq s\leq\e^u$ 已获证明，则 $\ell\leq\log s\leq u$。
若 $q\e^{-s}$ 已被一个非负有理区间包围，
利用 $z\mapsto\min\{1,c_1z,c_2z^2\}$ 的单调性，
即可包围 $g_q(\e^s)$。对差
$g_q(\e^{u_i})-g_q(\e^{u_{i+1}})$ 取上界时，
必须用前项的上端减后项的下端；
两个上端相减不构成合法上界。
生成的高度包围证明显式保留前后两个核区间，保证了这一方向。
对 $\exp$ 的有理区间由 Mathlib 的
\texttt{dyadic\_interval} 产生证明项，
而不是接受外部 Arb 计算的成功标记。

完整整数行表位于
\texttt{M2PilotIntegerData} 和 \texttt{M3FarIntegerData}；
逐行包围证明位于相应的
\texttt{Enclosures}、\texttt{FactorBounds} 模块。
例如，中段高度第 $0$ 行最终使用
\[
 G_0=6842520084211003,\qquad B_0=98272331280;
\]
每一行都按同一分母 $2^{48}$ 解释。
这一例子仅说明证书格式，最终不等式使用全部行。
正文给出了所有行的生成对象、核对条件和最终有限算式，
大型整数表则以可逐项读取的源码数据提供。

\subsection{中段试点：区间 [715/4,179]}

中段试点的参数完全固定为
\begin{equation}\label{cert:middle-parameters}
 \begin{gathered}
 A=\frac{715}{4},\quad C=179,\quad k=6,\quad
 q=1+\frac{\e^{179/6}}6,\quad J=1536,\quad M=988,\\
 \sigma_j=\frac{2048+j}{4096}\quad(0\leq j\leq1536),
 \qquad r(j)=\left\lfloor\frac j{128}\right\rfloor,\\
 u_i=\frac{1841+i}{64}\quad(0\leq i\leq988),\qquad
 U=\frac{2829}{64},\qquad
 p_0=\frac13,\quad v_0=\frac{2057}{64}.
 \end{gathered}
\end{equation}
每个左端点的行号为 $0,\ldots,11$，网格在
$1/2$ 开始并在 $7/8$ 结束。高度网格步长为 $1/64$；
在证书双重和中使用右端点 $s_i=u_{i+1}$。
因此不存在把 $988$ 个区间误解为 $988$ 个含首尾端点的节点。

为验证高度条件及 Taylor 区间，只需以下有限标量界：
\begin{gather}
 H\leq\e^{1841/64},\qquad
 q\e^{-1841/64}\geq\frac{31}{64},\label{cert:middle-start}\\
 \e^{215/64}\leq\frac{1842}{64},\qquad
 \frac{2829}{64}\leq\e^{284/64}. \label{cert:middle-log}
\end{gather}
第一个条件保证所有 CS 使用高度合法。第二个条件保证
$c_1q\e^{-u_0}\geq1$，且
$c_2q^2\e^{-2u_0}\geq31>1$，故 $g_q(\e^{u_0})=1$。
式 \eqref{cert:middle-log} 给出
$215/64\leq\log s_i\leq284/64$，于是
\[
 0\leq a_j\leq\frac23,\qquad
 0\leq b_i=s_i+\log s_i-\frac{2057}{64}\leq\frac{33}{2}.
\]
特别地 $a_jb_i\leq11$。这些界分别在
\leanref{M2PilotConstants}{start_logKernel_eq_one}
和 \leanref{M2PilotGrid}{sigma_height_product_bounds} 中落实。

对全部 $1536+988$ 行验证 \eqref{cert:dyadic-bounds} 后，
整数矩算式 \eqref{cert:rational} 满足
\begin{equation}\label{cert:middle-main}
 \mathcal M\leq\mathcal R\leq\frac{3759}{4000}=0.93975.
\end{equation}
实数因子与整数表的连接由
\leanref{M2PilotFactorBounds}{sigmaDyadicBounds}、
\leanref{M2PilotFactorBounds}{heightDyadicBounds} 给出；
有限有理不等式为
\leanref{M2PilotIntegerCertificate}{rationalLowRank20_le}。
该不等式使用普通内核归约
\texttt{decide +kernel}，并非将一个外部数值测试布尔值假定为真。

\subsection{远段试点：区间 [400,1601/4]}

远段使用同一个梯形权和同一个预算结构，但采用四次幂短区间长度：
\begin{equation}\label{cert:far-parameters}
 \begin{gathered}
 A=400,\quad C=\frac{1601}{4},\quad k=4,\quad
 q=1+\frac{\e^{(1601/4)/4}}4,\quad J=384,\quad M=125,\\
 \sigma_j=\frac{512+j}{1024}\quad(0\leq j\leq384),
 \qquad r(j)=\left\lfloor\frac j{32}\right\rfloor,\\
 u_i=\frac{795+i}{8}\quad(0\leq i\leq125),\qquad
 U=115,\qquad p_0=\frac13,\quad v_0=104.
 \end{gathered}
\end{equation}
所需标量界为
\[
 H\leq\e^{795/8},\qquad q\e^{-795/8}\geq\frac{31}{64},
 \qquad \e^{9/2}\leq\frac{796}{8},\qquad115\leq\e^5.
\]
由此 $0\leq a_j\leq2/3$、$0\leq b_i\leq16$，
所以 $a_jb_i\leq32/3<11$。
在相同的 $48$ 位有理格式下，完整的 $384+125$ 行证书给出
\begin{equation}\label{cert:far-main}
 \mathcal M\leq\mathcal R\leq\frac9{20}=0.45.
\end{equation}
对应 \leanref{M3FarGrid}{sigma_height_product_bounds}
和 \leanref{M3FarIntegerCertificate}{rationalLowRank20_le}。
远段证书没有使用历史稿的多项式权；当前远段的全部主张均以
本节的梯形核和式 \eqref{cert:far-parameters} 为准。

\subsection{辅助项、素数幂和平凡零点均已计入}

两个试点中，$B_{2,r}\leq10$、$B_{3,r}\leq168$。
实部权望远镜相消给出
\[
 \sum_jW_j(A)=\e^{-A/8}-\e^{-A/2}\leq\e^{-A/8},
 \quad D_A\leq10\e^{-A/8},\quad E_A\leq168\e^{-A/8}.
\]
高度质量也望远镜相消：
$\sum_i\mu_i=1-g_q(\e^U)\leq1$。
在两个试点中，最后一项
$32c_2q^2\e^{-2U}U^2$ 都不超过 $1$。
所以中段 $\sum_i d_i\leq2(89/2)^2+1<4000$，
远段 $\sum_i d_i\leq2\cdot115^2+1<30000$。
又因 $s_i\geq28$，有
\[
 D_A+\frac{E_A}{s_i}+\frac{\e^{-A/2+s_i}}{s_i}
 \leq16\e^{-A/8}+\frac{\e^{-A/2+U}}{28}.
\]
于是两个辅助预算分别不超过
\begin{align}
 \mathcal A_{\text{中}}
 &\leq4000\left(16\e^{-715/32}+\frac{\e^{-2891/64}}{28}\right)
       \leq\frac1{50000},\label{cert:middle-aux}\\
 \mathcal A_{\text{远}}
 &\leq30000\left(16\e^{-50}+\frac{\e^{-85}}{28}\right)
       \leq\frac1{50000}. \label{cert:far-aux}
\end{align}
最后两个都是有理端点上的指数不等式，另有区间证明，
不是从主项证书的剩余小数位中猜出的误差。
对应 \leanref{M2PilotAuxiliary}{auxiliaryBudget_le_small}
和 \leanref{M3FarAuxiliary}{auxiliaryBudget_le_small}。

由 \eqref{tr:Pcell}，在两个试点各自的 $[A,C]$ 上均有
\begin{equation}\label{cert:smallPV}
 P_k(A,C)\leq\frac1{10^6},\qquad V_A\leq\frac1{10^6}.
\end{equation}
这些估计留有很大余量。例如中段可用
$\epsilon_6(A)\leq1$、$\sqrt{1+\epsilon_6(A)}\leq2$、
$\log(1+\epsilon_6(A))\leq1$、$\log2\geq1/2$，
以及 $\e^{-A/3}\leq\e^{-50}\leq10^{-11}$。
远段的 $\e^{-A/4}=\e^{-100}$ 更小。
平凡零点则直接用 $V_A=\e^{-3A}/(1-\e^{-2A})$。
证明中并未将这些项设成零。

\begin{proposition}[两个完整试点证书]\label{cert:pilots}
在 QRH$(7/8)$、CS v2、Rosser 零点计数和积分显式公式四个输入下，
有以下结论：
\begin{align*}
 \frac{715}{4}\leq\log x\leq179
 &\Longrightarrow \exists p\text{ 为素数},\
                 x<p<x+6x^{5/6},\\
 400\leq\log x\leq\frac{1601}{4}
 &\Longrightarrow \exists p\text{ 为素数},\
                 x<p<x+4x^{3/4}.
\end{align*}
这里 $x>0$，所有区间端点均包含在适用域内。
\end{proposition}
\begin{proof}
中段的完整归一化预算为
\[
 Z+P+V\leq\frac{3759}{4000}+\frac1{50000}
                          +\frac2{10^6}
 =\frac{234943}{250000}=0.939772<\frac{47}{50}<1.
\]
远段的完整预算为
\[
 Z+P+V\leq\frac9{20}+\frac1{50000}+\frac2{10^6}
 =\frac{225011}{500000}=0.450022<\frac12<1.
\]
分别使用命题 \ref{dens:zero-budget} 和命题
\ref{tr:prime-criterion} 即得。
\end{proof}
对应 \leanref{M2PilotResult}{pilot_short_interval_from_inputs}
和 \leanref{M3FarResult}{pilot_short_interval_from_inputs}。
有限高度 RH 尚未用于这两个试点；$H$ 在这里仅是 CS 密度估计的
适用高度。下一节才在向下覆盖中段时使用有限 RH。
